\documentclass[11pt,a4paper]{article}

% --- arXiv-style preamble -----------------------------------------------------
\usepackage[utf8]{inputenc}
\usepackage[T1]{fontenc}
\usepackage{lmodern}
\usepackage[margin=1in]{geometry}
\usepackage{amsmath,amssymb,amsthm}
\usepackage{graphicx}
\usepackage{xcolor}
\usepackage[round,authoryear]{natbib}
\usepackage{booktabs}
\usepackage[hidelinks,colorlinks=true,linkcolor=blue,citecolor=blue,urlcolor=blue]{hyperref}

\bibliographystyle{plainnat}

% --- Title and authors --------------------------------------------------------
\title{Oracle3: An open-source autonomous trading agent for prediction markets
       with a Wang Transform pricing engine}

\author{%
  Yicheng Yang\thanks{Corresponding author: \href{mailto:ethanyang85@outlook.com}{ethanyang85@outlook.com}. ORCID: \href{https://orcid.org/0009-0000-7973-6931}{0009-0000-7973-6931}.}\\
  University of Illinois Urbana--Champaign, USA
}

\date{7 May 2026}

% --- Document ------------------------------------------------------------------
\begin{document}
\maketitle

\begin{abstract}
Prediction markets price binary contracts at systematically biased levels: a
contract whose objective probability is 50\% trades on average around 57 cents,
an empirical regularity known as the favorite--longshot bias
\citep{thaler:1988}. Despite a long literature documenting this distortion,
most open-source trading bots ignore it, treating market prices as unbiased
estimates of probability. \texttt{Oracle3} is a Python framework that
operationalizes a peer-reviewed risk-neutral pricing model for binary outcome
contracts and uses it to drive automated trading across multiple venues
(Kalshi, Polymarket, and Solana-based DFlow). At its core is a Wang Transform
\citep{wang:2000,wang:2002} calibrated by maximum-likelihood estimation on
291{,}309 resolved contracts spanning six platforms, with hierarchical
covariates for volume, days-to-expiry, and contract moneyness. The library
exposes the model as a fair-value engine and pairs it with eight
constraint-based arbitrage strategies, statistical-arbitrage strategies,
model-Greek-driven sizing, and a risk manager --- all wired into an
event-driven async trading core with snapshot persistence, killswitch support,
and on-chain audit trails. Source code is available at
\url{https://github.com/YichengYang-Ethan/oracle3} under the Apache~2.0 license,
with archived release at Zenodo (\href{https://doi.org/10.5281/zenodo.20062549}{10.5281/zenodo.20062549}).
\end{abstract}

\section{Introduction}
\label{sec:intro}

Quantitative research on prediction-market pricing has accelerated, but the gap
between empirical pricing literature and reproducible, deployable software
remains wide. Existing open-source bots (e.g., \texttt{polymarket-py},
\texttt{kalshi-python}) are thin client wrappers; they do not estimate fair
value, do not operationalize known pricing distortions, and offer no shared
abstraction across venues. Conversely, academic replication packages typically
stop at calibration and rarely produce a runnable trading agent.
\texttt{Oracle3} closes this gap by shipping the maximum-likelihood estimates
from \citet{yang:2026} directly as a real-time pricing engine, exposing the
same hierarchical covariate model used in the paper, and connecting it to live
order-book data, execution, and risk management.

The intended audience is twofold. For researchers studying prediction-market
efficiency, the favorite--longshot bias \citep{thaler:1988,snowberg:2010}, and
risk-neutral pricing \citep{harrison:1979,wang:2000}, \texttt{Oracle3}
provides a reproducible, well-tested implementation of a multi-venue
pricing-and-trading pipeline that can be used as a baseline for new strategies
or as a calibration testbed. For practitioners, it provides a transparent,
auditable agent whose every signal is grounded in an explicit pricing model
rather than opaque heuristics.

\paragraph{State of the field.}
Several Python tools touch the prediction-market space, but each is partial.
\texttt{py-clob-client} and \texttt{kalshi-python} are vendor SDKs without
modeling. The \texttt{crypto-trading-bot} family targets centralized exchanges
and currency markets, not binary outcome markets. Research code accompanying
empirical-finance papers is typically Jupyter-driven, focused on calibration
rather than deployment. Outside Python, decentralized-finance projects on
Solana such as \texttt{dflow-program} provide on-chain primitives but no
off-chain modeling layer. \texttt{Oracle3} is, to our knowledge, the first
open-source library that combines (i) calibrated probabilistic pricing for
binary contracts, (ii) constraint-based arbitrage detection across
heterogeneous venues, and (iii) production-grade execution with on-chain
settlement.

\section{Methodology}
\label{sec:method}

The core pricing primitive is the Wang Transform applied to binary contracts:
%
\begin{equation}
  p^{\mathrm{mkt}} \;=\; \Phi\!\left(\Phi^{-1}(p^{*}) + \lambda\right),
  \label{eq:wang}
\end{equation}
%
where $p^{*}$ is the objective (true) probability of the binary outcome,
$p^{\mathrm{mkt}}$ is the market-clearing price, $\Phi$ is the standard normal
CDF, and $\lambda$ is a market-price-of-risk parameter that absorbs systematic
risk preferences and microstructure frictions. A positive $\lambda$ shifts
prices toward 0.5 from below and toward 1 from above, capturing both the
favorite--longshot bias documented by \citet{thaler:1988,snowberg:2010} and the
risk-neutralization argument of \citet{wang:2000,wang:2002}, which generalizes
the no-arbitrage pricing of \citet{harrison:1979} to incomplete markets.

\paragraph{Calibration.}
We estimate $\lambda$ by maximum likelihood on a panel of $N = 291{,}309$
resolved contracts collected from six platforms. The pooled estimate is
$\hat\lambda = 0.183$. To capture cross-sectional and lifecycle heterogeneity,
we additionally fit a hierarchical covariate specification:
%
\begin{equation}
  \lambda_i \;=\; 0.259 \;-\; 0.072\,\ln(1+V_i) \;+\; 0.143\,\ln(1+D_i)
                    \;-\; 0.477\,\lvert p_i - 0.5 \rvert,
  \label{eq:lambda}
\end{equation}
%
where $V_i$ is daily traded volume, $D_i$ is days remaining to expiry, and
$\lvert p_i - 0.5 \rvert$ is contract moneyness. All three covariates have
signs consistent with the favorite--longshot literature
\citep{thaler:1988,snowberg:2010,wolfers:2004,manski:2006,ottaviani:2007}: less
liquid, longer-dated, and far-out-of-the-money contracts exhibit larger risk
premia. The coefficient estimates and full replication procedure appear in the
working paper \citep{yang:2026}.

\paragraph{Streaming recalibration.}
A production trading agent must adapt to regime change without forgetting
historical structure. \texttt{Oracle3}'s online recalibrator combines the
batch MLE estimate with an exponentially weighted moving-average (EWMA)
estimator over recent resolutions, and applies category-level shrinkage so
sparse subgroups (e.g., a freshly listed event class) inherit prior pooled
$\lambda$ until enough data accumulates. This bridges the academic batch
calibration of \citet{yang:2026} with a real-time deployment.

\paragraph{Sizing and risk.}
Position sizing follows a fractional Kelly criterion \citep{kelly:1956}
applied to the model-implied edge $p^{*} - p^{\mathrm{mkt}}$, capped by a risk
budget. Multi-leg arbitrage trades are sized so that worst-case loss across
all legs respects the agent's drawdown limits.

\section{Implementation}
\label{sec:impl}

\texttt{Oracle3} is organized as four layers.

\paragraph{1. Pricing engine (\texttt{oracle3/pricing/}).}
Implements equations \eqref{eq:wang} and \eqref{eq:lambda} as a vectorized
NumPy/PyTorch module, with the global $\hat\lambda = 0.183$ available as a
fast path and the hierarchical covariate model used when contract metadata
are present. The online recalibrator combines batch MLE with a streaming EWMA
estimator and category shrinkage, so the model adapts to regime change
without forgetting prior data.

\paragraph{2. Strategy layer (\texttt{oracle3/strategy/}).}
Eight constraint-based strategies each enforce a probability-axiom invariant
(cross-market identity, exclusivity bound $P(A) + P(B) \le 1$, implication
monotonicity, conditional bounds, event-sum unity, and structural
relationships derived from the calibrated model) and emit signals when
violations exceed configurable thresholds. Two model-driven strategies trade
fair-value divergence and a ``premium decay'' lifecycle effect predicted by
the $D$-coefficient in equation \eqref{eq:lambda}. Statistical-arbitrage
strategies (cointegration spread, lead-lag) provide complementary signals.

\paragraph{3. Trading core (\texttt{oracle3/core/}, \texttt{oracle3/trader/}).}
An async event loop dispatches strategy signals to a \texttt{SpreadExecutor}
that posts multi-leg orders atomically and unwinds via LIFO on partial fills,
eliminating naked legs. A dual-layer risk manager enforces position,
drawdown, and exposure limits locally, and additionally calls the Solana
\texttt{simulateTransaction} RPC for pre-flight on-chain checks before any
DFlow submission.

\paragraph{4. Infrastructure (\texttt{oracle3/cli/}, \texttt{oracle3/data/}).}
A Click-based CLI exposes market discovery, live dashboards, backtests, and
paper-run modes. Snapshot persistence and a Unix-socket control plane enable
pause/resume/killswitch operations on long-running agents. Solana submissions
go through Jito bundles, with a Memo-program audit trail for every fill.

\paragraph{Quality and reproducibility.}
The codebase is type-checked with \texttt{mypy}, linted with \texttt{ruff},
and validated by 633 unit and integration tests on every push, enforced by
GitHub Actions. The design has been used to run paper-trading sessions on
Polymarket and Kalshi, and the pricing engine is shared with the replication
package of \citet{yang:2026}.

\section{Conclusion}
\label{sec:conclusion}

\texttt{Oracle3} narrows the gap between empirical prediction-market pricing
research and deployable open-source trading software. By packaging a
peer-reviewed Wang Transform calibration as a real-time engine, wiring it
into multi-venue execution with on-chain settlement, and exposing every
signal through an auditable strategy layer, the library provides a common
substrate for both empirical research and live experimentation in
prediction-market microstructure.

\section*{Acknowledgements}

The author thanks the maintainers of the Polymarket, Kalshi, and DFlow APIs
for documentation; the prediction-market research community whose empirical
work motivates this software; and contributors who reported issues and
suggested strategies. Any errors are the author's own.

\section*{Code and data availability}

The complete source code is available at
\url{https://github.com/YichengYang-Ethan/oracle3} under the Apache~2.0
license. An archived release with persistent identifier is deposited at
Zenodo, DOI \href{https://doi.org/10.5281/zenodo.20062549}{10.5281/zenodo.20062549}.
The companion empirical paper is available as an SSRN working paper:
\url{https://papers.ssrn.com/sol3/papers.cfm?abstract_id=6468338}.

\bibliography{paper}

\end{document}
